\documentclass[UTF8,zihao=-4]{ctexart}
\usepackage[a4paper,margin=2.1cm]{geometry}
\usepackage{hyperref}
\usepackage{array}
\usepackage{longtable}
\usepackage{verbatim}
\hypersetup{hidelinks}
\setlength{\parindent}{0pt}
\setlength{\parskip}{0.55em}
\emergencystretch=4em
\sloppy
\title{期末试题答案汇总}
\author{}
\date{}
\begin{document}
\maketitle
\setcounter{tocdepth}{1}
\tableofcontents
\newpage
\section{编译原理期末试题答案汇总}

\subsection{使用结构}

\begin{enumerate}
\item 先看“题型速查”，按同类题合并复习。
\item 再看“通用答题模板”，掌握自动机、LL、LR、回填、SSA、调度、符号执行、指针分析。
\item 最后按年份看题目对应答案。
\end{enumerate}

\subsection{题型速查}

\begin{itemize}
\item 正则表达式、Thompson NFA、子集构造 DFA、DFA 最小化：2021、2022、2023、2025。复习重点是逐层画 NFA，DFA 状态写成 NFA 状态集合，最小化按分组细化。
\item LL(1)：2021、2022、2023、2025。复习重点是消除左递归、FIRST、FOLLOW、预测分析表、冲突判断。
\item LR/SLR/LR(1)：2021、2022、2023。复习重点是 closure、goto、项目集、归约条件、冲突说明。
\item ANTLR4 优先级上升：2022、2023。复习重点是运算符优先级、结合性、改写文法、分析树。
\item 语法制导定义与翻译：2022、2023。复习重点是属性定义、综合/继承属性、语义规则。
\item 中间代码与回填：2022、2023。复习重点是 truelist、falselist、nextlist、while/if 回填。
\item SSA 与支配关系：2025。复习重点是 Dom、Post-Dom、Dom Frontier、支配树。
\item 指令调度：2025。复习重点是 True/Anti/Output 依赖、依赖图、重命名、拓扑调度。
\item 符号执行与 SAT：2025。复习重点是敏感性、CNF、Tseitin、同时可满足。
\item 指针别名：2025。复习重点是 Andersen 约束、图闭包、points-to 集合。
\end{itemize}

\subsection{通用答题模板}

\subsubsection{模板 1：Thompson 构造}

\begin{quote}
\footnotesize
\begin{verbatim}
epsilon:  start --epsilon--> end
a      :  start --a--> end
A|B    :  新 start epsilon 到 A.start/B.start，A.end/B.end epsilon 到新 end
AB     :  A.end epsilon 到 B.start
A*     :  新 start epsilon 到 A.start 和新 end；A.end epsilon 到 A.start 和新 end
\end{verbatim}
\end{quote}

\subsubsection{模板 2：子集构造}

\begin{quote}
\footnotesize
\begin{verbatim}
Start = epsilon-closure({NFA start})
对每个 DFA 状态 S 和字符 a:
  T = epsilon-closure(move(S, a))
  若 T 新出现，加入状态集合
终止状态：S cap F_NFA != emptyset
\end{verbatim}
\end{quote}

\subsubsection{模板 3：DFA 最小化}

\begin{quote}
\footnotesize
\begin{verbatim}
P0 = {F, Q-F}
按每个输入字符的转移目标所在组拆分
重复直到分组不再变化
每个分组合并成一个状态
\end{verbatim}
\end{quote}

\subsubsection{模板 4：LL(1)}

\begin{quote}
\footnotesize
\begin{verbatim}
1. 消除左递归
2. 计算 FIRST
3. 计算 FOLLOW
4. 对 A -> alpha:
   FIRST(alpha)-{epsilon} 中每个 a 填 M[A,a]
   若 epsilon in FIRST(alpha)，FOLLOW(A) 中每个 b 填 M[A,b]
5. 表中无多重入口，则是 LL(1)
\end{verbatim}
\end{quote}

\subsubsection{模板 5：LR(1)}

\begin{quote}
\footnotesize
\begin{verbatim}
项目: [A -> alpha . beta, a]
closure:
  若 [A -> alpha . B beta, a] 在 I 中，
  对 B -> gamma 和 b in FIRST(beta a)，加入 [B -> . gamma, b]
goto(I,X):
  closure({[A -> alpha X . beta, a]})
归约:
  [A -> alpha ., a] 只在 a 列填 reduce
\end{verbatim}
\end{quote}

\subsubsection{模板 6：回填}

\begin{quote}
\footnotesize
\begin{verbatim}
B.truelist, B.falselist, S.nextlist
makelist(i), merge(p,q), backpatch(p,L)
while:
  记录循环入口 M1
  计算 B
  回填 B.truelist 到循环体入口
  回填循环体 nextlist 到 M1
  发出 goto M1
  S.nextlist = B.falselist
\end{verbatim}
\end{quote}

\subsubsection{模板 7：支配与支配边界}

\begin{quote}
\footnotesize
\begin{verbatim}
Dom(A): 从入口到 A 的所有路径都经过的块
PostDom(A): 从 A 到出口的所有路径都经过的块
idom(A): A 最近的严格支配者
DF(B): B 支配某个前驱，但不严格支配该节点的位置
\end{verbatim}
\end{quote}

\subsubsection{模板 8：Andersen 约束}

\begin{quote}
\footnotesize
\begin{verbatim}
y = &x : x in pts(y)
y = x  : pts(x) subset pts(y)
*y = x : for v in pts(y), pts(x) subset pts(v)
y = *x : for v in pts(x), pts(v) subset pts(y)
\end{verbatim}
\end{quote}

\subsection{2025 答案}

\subsubsection{一、词法分析：正则表达式到 NFA、DFA、最小化}

题型：把题面正则表达式按 Thompson 构造为 NFA，再做子集构造和 DFA 最小化。

答案步骤：

\begin{enumerate}
\item 按括号最内层到最外层拆正则表达式。
\item 对 \texttt{epsilon}、字符、连接、选择、闭包分别套 Thompson 模板。
\item 标号 NFA 状态，写出 epsilon-closure。
\item 从起始集合开始列 DFA 转移表。
\item 含 NFA 终止状态的集合标为 DFA 终止状态。
\item 按终止/非终止初分组，迭代拆分，得到最小 DFA。
\end{enumerate}

\subsubsection{二、语法分析：消除左递归、FIRST/FOLLOW、预测表}

题面文法：

\begin{quote}
\footnotesize
\begin{verbatim}
A -> B C
B -> B P | r
C -> t | u
\end{verbatim}
\end{quote}

消除 \texttt{B} 的直接左递归：

\begin{quote}
\footnotesize
\begin{verbatim}
A  -> B C
B  -> r B'
B' -> P B' | epsilon
C  -> t | u
\end{verbatim}
\end{quote}

FIRST：

\begin{quote}
\footnotesize
\begin{verbatim}
FIRST(B)  = {r}
FIRST(B') = {P, epsilon}
FIRST(C)  = {t, u}
FIRST(A)  = {r}
\end{verbatim}
\end{quote}

FOLLOW：

\begin{quote}
\footnotesize
\begin{verbatim}
FOLLOW(A)  = {$}
FOLLOW(B)  = {t, u}
FOLLOW(B') = {t, u}
FOLLOW(C)  = {$}
\end{verbatim}
\end{quote}

预测表：

\begin{quote}
\footnotesize
\begin{verbatim}
M[A, r] = A -> B C
M[B, r] = B -> r B'
M[B', P] = B' -> P B'
M[B', t] = B' -> epsilon
M[B', u] = B' -> epsilon
M[C, t] = C -> t
M[C, u] = C -> u
\end{verbatim}
\end{quote}

结论：表中无多重入口，消除左递归后的文法是 LL(1)。

\subsubsection{三、指令调度}

目的：

\begin{itemize}
\item 让独立指令靠近，减少流水线停顿。
\item 提高硬件并行执行利用率。
\end{itemize}

依赖定义：

\begin{quote}
\footnotesize
\begin{verbatim}
True dependence / RAW  : i 写 x，j 读 x，i 必须在 j 前。
Anti dependence / WAR  : i 读 x，j 写 x，j 提前会破坏 i 读到的值。
Output dependence / WAW: i 写 x，j 写 x，顺序影响最终写入值。
\end{verbatim}
\end{quote}

构建依赖图：

\begin{enumerate}
\item 每条指令一个节点。
\item RAW 画实线边。
\item WAR 和 WAW 画虚线边。
\item 寄存器重命名可消除 WAR 和 WAW。
\item 对保留依赖图做拓扑排序，得到调度序列。
\end{enumerate}

\subsubsection{四、支配问题}

定义：

\begin{quote}
\footnotesize
\begin{verbatim}
Dom(n): 从入口到 n 的每条路径都经过的节点集合。
PostDom(n): 从 n 到出口的每条路径都经过的节点集合。
Dom Frontier(n): n 支配某个前驱，但不严格支配该节点的节点集合。
\end{verbatim}
\end{quote}

支配关系图是森林的证明：

\begin{enumerate}
\item 除入口节点外，每个节点有唯一直接支配者。
\item 直接支配者严格支配该节点。
\item 严格支配关系无环。
\item 因此每个节点至多一个父节点，且无环，构成以入口为根的树；多入口情形构成森林。
\end{enumerate}

\subsubsection{五、符号执行}

敏感性定义：

\begin{quote}
\footnotesize
\begin{verbatim}
Path-sensitive: 区分不同执行路径。
Flow-sensitive: 区分程序点和语句顺序。
Context-sensitive: 区分函数调用上下文。
\end{verbatim}
\end{quote}

CNF：

\begin{quote}
\footnotesize
\begin{verbatim}
若干子句的合取，每个子句是文字的析取。
(l11 or l12) and (l21 or l22 or l23)
\end{verbatim}
\end{quote}

Tseitin 变换：

\begin{enumerate}
\item 为每个子表达式建新变量。
\item 加入新变量与子表达式等价的 CNF 约束。
\item 加入根变量为真。
\item 得到与原式同时可满足的 CNF。
\end{enumerate}

题中同时可满足证明模板：

\begin{itemize}
\item 对辅助变量 \texttt{c} 分情况。
\item 若左式可满足，则赋值满足其中一个含 \texttt{c} 子句和另一个含 \texttt{not c} 子句，去掉 \texttt{c} 后得到大析取式可满足。
\item 若大析取式可满足，则若某个 \texttt{a\_i} 为真，令 \texttt{c=false}；若某个 \texttt{b\_j} 为真，令 \texttt{c=true}，左式可满足。
\item 两边可满足性等价。
\end{itemize}

\subsubsection{六、指针别名：Andersen}

题面语句模板：

\begin{quote}
\footnotesize
\begin{verbatim}
p = &a;
q = &b;
*p = q;
r = &c;
s = p;
t = *p;
*s = r;
\end{verbatim}
\end{quote}

约束：

\begin{quote}
\footnotesize
\begin{verbatim}
a in pts(p)
b in pts(q)
for v in pts(p), pts(q) subset pts(v)
c in pts(r)
pts(p) subset pts(s)
for v in pts(p), pts(v) subset pts(t)
for v in pts(s), pts(r) subset pts(v)
\end{verbatim}
\end{quote}

闭包推导：

\begin{quote}
\footnotesize
\begin{verbatim}
pts(p) = {a}
pts(q) = {b}
由 *p=q 得 pts(a) 包含 {b}
pts(r) = {c}
由 s=p 得 pts(s) 包含 {a}
由 t=*p 且 pts(p)={a} 得 pts(a) subset pts(t)，所以 pts(t) 包含 {b}
由 *s=r 且 pts(s)={a} 得 pts(a) 包含 {c}
最终 pts(a) = {b,c}，pts(t) 包含 {b,c}
\end{verbatim}
\end{quote}

\subsection{2023 答案}

\subsubsection{题 1：正则表达式与自动机}

题型同 2025 第一题。答案按 Thompson、子集构造、DFA 最小化三步写。

\subsubsection{题 2：LL 语法分析}

文法：

\begin{quote}
\footnotesize
\begin{verbatim}
S -> cS | AB
A -> aA | epsilon
B -> bBa | d
\end{verbatim}
\end{quote}

FIRST：

\begin{quote}
\footnotesize
\begin{verbatim}
FIRST(A) = {a, epsilon}
FIRST(B) = {b, d}
FIRST(S) = {c, a, b, d}
\end{verbatim}
\end{quote}

FOLLOW：

\begin{quote}
\footnotesize
\begin{verbatim}
FOLLOW(S) = {$}
FOLLOW(A) = FIRST(B) = {b, d}
FOLLOW(B) = FOLLOW(S) = {$}
\end{verbatim}
\end{quote}

预测表：

\begin{quote}
\footnotesize
\begin{verbatim}
M[S,c] = S -> cS
M[S,a] = S -> AB
M[S,b] = S -> AB
M[S,d] = S -> AB
M[A,a] = A -> aA
M[A,b] = A -> epsilon
M[A,d] = A -> epsilon
M[B,b] = B -> bBa
M[B,d] = B -> d
\end{verbatim}
\end{quote}

结论：无冲突，是 LL(1) 文法。

\subsubsection{题 3：ANTLR4 优先级上升}

原文法：

\begin{quote}
\footnotesize
\begin{verbatim}
S -> E
E -> E ! | E ^ E | E + E | INT
\end{verbatim}
\end{quote}

优先级：

\begin{quote}
\footnotesize
\begin{verbatim}
! 最高，后缀
^ 右结合
+ 最低，左结合
\end{verbatim}
\end{quote}

答题要点：

\begin{itemize}
\item 先把 \texttt{5!!} 归为后缀链。
\item \texttt{2\textasciicircum{}3\textasciicircum{}4} 按右结合归为 \texttt{2\textasciicircum{}(3\textasciicircum{}4)}。
\item 最外层用 \texttt{+} 左结合：\texttt{(1 + (2\textasciicircum{}(3\textasciicircum{}4))) + ((5!)!)}。
\end{itemize}

\subsubsection{题 4：语法制导翻译}

任务：二进制小数转十进制。

可定义属性：

\begin{quote}
\footnotesize
\begin{verbatim}
S.val
L.val, L.len
B.val
\end{verbatim}
\end{quote}

规则：

\begin{quote}
\footnotesize
\begin{verbatim}
S -> L1 . L2   S.val = L1.val + L2.val / 2^(L2.len)
S -> L         S.val = L.val
L -> L1 B      L.val = L1.val * 2 + B.val; L.len = L1.len + 1
L -> B         L.val = B.val; L.len = 1
B -> 0         B.val = 0
B -> 1         B.val = 1
\end{verbatim}
\end{quote}

\subsubsection{题 5：中间代码生成}

二分查找 while 部分按回填模板写。关键结构：

\begin{quote}
\footnotesize
\begin{verbatim}
Lbegin:
if low <= high goto Lbody
goto Lexit
Lbody:
mid = (low + high) / 2
if dict[mid] < key goto Lless
goto Lnotless
Lless:
low = mid + 1
goto Lbegin
Lnotless:
if dict[mid] > key goto Lgreater
goto Lequal
Lgreater:
high = mid - 1
goto Lbegin
Lequal:
return mid
Lexit:
return -1
\end{verbatim}
\end{quote}

\subsection{2022 答案}

\subsubsection{题 1：正则表达式与自动机}

表达式：\texttt{(a | b)* a+}。

作答：

\begin{itemize}
\item 对 \texttt{(a|b)} 构造选择 NFA。
\item 对 \texttt{(a|b)*} 加星闭包。
\item 对 \texttt{a+} 写成 \texttt{a a*}。
\item 将两部分连接。
\item 子集构造时 DFA 起始状态为起始点 epsilon-closure。
\item 终止状态为包含 NFA 终点的集合。
\end{itemize}

\subsubsection{题 2：LL 语法分析}

文法：

\begin{quote}
\footnotesize
\begin{verbatim}
D -> T L
T -> int | real
L -> L , id | id
\end{verbatim}
\end{quote}

消除 \texttt{L} 左递归：

\begin{quote}
\footnotesize
\begin{verbatim}
D  -> T L
T  -> int | real
L  -> id L'
L' -> , id L' | epsilon
\end{verbatim}
\end{quote}

FIRST：

\begin{quote}
\footnotesize
\begin{verbatim}
FIRST(D)  = {int, real}
FIRST(T)  = {int, real}
FIRST(L)  = {id}
FIRST(L') = {',', epsilon}
\end{verbatim}
\end{quote}

FOLLOW：

\begin{quote}
\footnotesize
\begin{verbatim}
FOLLOW(D)  = {$}
FOLLOW(T)  = {id}
FOLLOW(L)  = {$}
FOLLOW(L') = {$}
\end{verbatim}
\end{quote}

预测表：

\begin{quote}
\footnotesize
\begin{verbatim}
M[D,int] = D -> T L
M[D,real] = D -> T L
M[T,int] = T -> int
M[T,real] = T -> real
M[L,id] = L -> id L'
M[L',,] = L' -> , id L'
M[L',$] = L' -> epsilon
\end{verbatim}
\end{quote}

结论：无冲突，是 LL(1)。

\subsubsection{题 5：语法制导定义}

任务 1：统计声明个数。

\begin{quote}
\footnotesize
\begin{verbatim}
P -> D              P.count = D.count
D -> D1 ; D2        D.count = D1.count + D2.count
D -> T id           D.count = 1
D -> func id {D1}   D.count = 1 + D1.count
\end{verbatim}
\end{quote}

任务 2：打印变量嵌套深度。

\begin{quote}
\footnotesize
\begin{verbatim}
P -> D              D.depth = 0
D -> D1 ; D2        D1.depth = D.depth; D2.depth = D.depth
D -> T id           print(id, D.depth)
D -> func id {D1}   D1.depth = D.depth + 1
\end{verbatim}
\end{quote}

\subsubsection{题 6：中间代码生成}

GCD while 部分：

\begin{quote}
\footnotesize
\begin{verbatim}
Lbegin:
if a != b goto Lbody
goto Lexit
Lbody:
if a > b goto Lthen
goto Lelse
Lthen:
a = a - b
goto Lbegin
Lelse:
b = b - a
goto Lbegin
Lexit:
return a
\end{verbatim}
\end{quote}

\subsection{2021 与 2020 答案}

PDF 文本抽取保留了题型主干，答案按通用模板作答：

\begin{itemize}
\item 正则表达式与自动机：Thompson NFA、子集构造 DFA、DFA 最小化。
\item LL：消除左递归、FIRST、FOLLOW、预测表、LL(1) 判断。
\item LR：FIRST/FOLLOW、LR(1) 自动机、LR(1) 分析表、冲突判断。
\item CFG/CNF：按 Chomsky Normal Form 要求改写产生式。
\item 语法制导翻译：定义综合属性和继承属性，写产生式对应语义规则。
\item 目标代码生成：按题面给定指令集合填空，遵守调用约定、条件跳转、栈帧和返回值规则。
\end{itemize}
\end{document}
